\section{An Optimal IR Resource Exchange Algorithm}
\label{sec:opt-non-ic-alg}

% Our first algorithm design goal will be to determine when there exists
% an optimal algorithm and to construct one in this case.
% \mycomment{I don't think the term policy is mentioned anywhere in durrett}
Our first goal is to determine whether an optimal algorithm exists, and to construct one when it does.
To do so,
we will draw on tools from martingale theory to design our exchange algorithm and, accordingly, adopt some of the 
terminology from this literature. In this context, the term \emph{policy} is used as a generalization of stopping 
times to partially ordered sets~\citep{mandelbaum1981optimal,williams1991probability}. We focus on a restricted class 
of such policies, which we term \emph{exchange policies}. These policies satisfy additional structural properties that 
enable the design of exchange algorithms (Definition~\ref{def:exchange-alg}). In Section~\ref{sec:exchangealgotopolicy}, 
we will first show how an exchange algorithm can be reformulated as an exchange policy.
Then, in Section~\ref{subsec:reduction-to-optimal-stopping}, we will show how we can use optimal stopping over the lattice to derive an optimal exchange algorithm.

\subsection{Reformulating an Exchange Algorithm as an Exchange Policy}
\label{sec:exchangealgotopolicy}

%\kkcomment{Can we not call this an 'exchange policy'? }




% We will use tools from martingale theory to design an exchange algorithm, and as such, borrow some verbiage from this literature.
% In this literature, the term 'policy' is used as a generalization of stopping times for partially ordered sets~\citep{mandelbaum1981optimal,williams1991probability,durrett2019probability}.
% We will consider a subset of such policies, which we call `exchange policies', which have additional properties and can be used to design exchange algorithms (Definition~\ref{def:exchange-alg}).
% We will now define these terms formally.

Suppose that an exchange 
\(
    \exchange
    =
    \rbr{
        \numdatafora,
        \numdataforb
    }
\)
has already occurred. In this case $\competitora$ has observed their value for the first 
$\numdatafora$ of $\competitorb$'s resources and 
$\competitorb$ has observed their value for the first $\numdataforb$ of $\competitora$'s resources.
To represent the information
available to an algorithm after $\competitora$ and $\competitorb$ report their values for the 
exchange 
\(
    \exchange=\rbr{
        \numdatafora, \numdataforb
    }\in\setexchanges
\), we define the $\sigma$-algebra
\begin{align*}
    \statefilts
    =
    \sigma\rbr{
        \cbr{
            \valaii{k}
        }_{k=0}^\numdatafora,
        \cbr{
            \valbjj{k}
        }_{k=0}^\numdataforb
    }
    \numberthis
    \label{eq:family-sig-alg}
\end{align*}  
and consider the family
\(
    \cbr{
        \statefilts
    }_{
        \exchange\in\setexchanges
    }
\).
Notice that for two exchanges
\(
    \exchangett{1},\exchangett{2}
    \in\setexchanges
\)
we have 
% \begin{align*}
\(
    \exchangett{1}\leq\exchangett{2}
    \iff 
    \statefiltss{
        \exchangett{1}
    }
    \subseteq
    \statefiltss{
        \exchangett{2}
    }
\).
% \end{align*}  
This reflects that larger exchanges contain more information about the agents' valuation functions.
Additionally, 
\(
    \statefiltss{0}
    =
    \cbr{\varnothing,\Omega}
\)
which captures that no information besides the distributions $\priora,\priorb$ of the valuation functions is known prior to 
$\competitora$ and $\competitorb$ exchanging any data.

From Defintion~\ref{def:exchange-alg}, we see that for any algorithm $\alg\in\classalgs$, there is a correspondence
between $\alg$ and the sequence of $\setexchanges$-valued random variables 
\(
    \rbr{
        \algstatett{0},
        \ldots
        \algstatett{\algendtime}
    }
\)
it induces. Moreover, it is the case that 
\(
    \algstatett{\timeind+1}
\)
is
\(
    \statefiltss{
        \algstatet
    }
\)-measurable, reflecting that at round $\timeind$, the choice of the next exchange
only depends on information revealed at the time of the current exchange.
Intuitively, an algorithm can be thought of as a way of traversing $\setexchanges\subset\NN^2$, starting at $0$, and moving up or right
each round,
and deciding when to stop.
We will shortly see that an algorithm is a special case of 
% what~\citet{mandelbaum1981optimal} defines to be a policy.
a policy, which is the generalization of a stopping time to a partially ordered set.
\begin{restatable}{definition}{policy}
    \label{def:policy} 
    \textup{
        (\citet{mandelbaum1981optimal})
    }
    Let $\osasstates$ be a countable partially ordered set. For $r,s\in S$ denote
    \(
        \sbr{r,\osasstate}
        =
        \cbr{
            u\in \osasstates: r\leq u\leq \osasstate
        } 
    \).
     Let 
    \(
        \osasdirectsuccessorss{
            \osasstate
        }
        =
        \cbr{
            u> \osasstate: \sbr{\osasstate,u}=\cbr{\osasstate,u}
        } 
    \) be the set of direct successors.
    Assume that $\sbr{r,\osasstate}$ is finite and there is a unique minimal element $0\in \osasstates$. 
    Let 
    \(
        \cbr{\statefilts}_{
            \osasstate\in \osasstates
        }
    \)
    be a family of $\sigma$-algebras where $\Fcal_r\subseteq\Fcal_\osasstate$ for $r\leq \osasstate$. A strategy is 
    a sequence of $\osasstates$-valued random variables $\cbr{\sigma_\timeind}_{\timeind\geq 0}$ such that 
    \(
        \sigma_\timeind \leq \sigma_{\timeind+1}
    \)
    and
    \(
        % \quad
        \sigma_{\timeind+1} \text{ is }
        \Fcal_{\sigma_\timeind}\text{-measurable}
    \).
    A policy $\pi$ is a pair
    \(
        \rbr{
            \cbr{
                \osasstrategyt
            }_{\timeind\geq 0},
            \osasstoptime
        }
    \)
    where $\osasstoptime$ is a stopping time with respect to the filtration
    \(
        \cbr{
            \timefiltt
        }_{\timeind\geq 0}
        =
        \cbr{
            \statefilt_{
                \strategyt
            }
        }_{\timeind\geq 0}
    \).
    Denote $\osaspolicies$ as the class of policies such that $\osasstrategytt{0}=0$.
\end{restatable}  

To limit the number of subscripts, we will occassionally write
$\stopstateforpolicy{\apolicy}$ as shorthand for 
$\strategytt{\stoptime}$.
In~\citet{mandelbaum1981optimal}, each state is associated with a reward $\osasrewards$, and the expected 
reward of a policy 
\(
    \rbr{
        \cbr{
            \sigma_t
        }_{t\geq 0},
        \tau
    }
\)
is defined as
\(
    \EE\sbr{
        \osasrewardss{
            \osasstrategytt{
                \osasstoptime
            }
        }
    }
\).
Because of this, when specifying a policy 
\(
    \rbr{
        \cbr{
            \osasstrategyt
        }_{\timeind\geq 0},
        \osasstoptime 
    }
    % \in
    % \osaspolicies
\),
we only need to worry about defining $\strategyt$ for $\timeind\leq\stoptime$.
In our setting, an algorithm can be viewed as a policy that satisfies additional structure. We refer to such policies as \emph{exchange policies} and define them formally below.


% For our problem, an algorithm can be viewed as a policy with two additional requirements, taking $S=\setexchanges$
% (equiped with the standard partially ordering over $\NN^2$), 
% the family of $\sigma$-algebras as given in~\eqref{eq:family-sig-alg},
% and the reward $\osasrewards$ associated to $\exchange\in\setexchanges$ to be $\rewards$. 
% % Thus we have the following equivalent definition of a resource exchange algorithm.
% This is formalized in the following definition.
% \kkcomment{write more clearly. this sounds like there are 3 additional requirements.}


\begin{restatable}{definition}{exchange-policy}
    \label{def:exchange-policy} 
    A resource exchange policy
    % $\alg$
    is a policy 
    \(
        \osaspolicy=
        \rbr{
            \cbr{
                \strategyt
            }_{\timeind\geq 0},
            \stoptime 
        }
        \in
        \osaspolicies
    \),
    as specified in Definition~\ref{def:policy}, taking
    \(
        \osasstates=\setexchanges
    \)
    (which has the unique minimal element $(0,0)$)
    equiped with the family of $\sigma$-algebras 
    \(
        \cbr{
            \statefilts
        }_{\exchange\in\setexchanges}
    \)
    given in~\eqref{eq:family-sig-alg}. We additionally require that
    \(
        \strategyt < \strategytt{\timeind+1}
    \)
    for
    \(
        \strategyt < \rbr{
            \numdataa,\numdatab
        },
    \)
    and that $\stoptime$ is finite.
    Denote $\setexchangepolicies$ as the class of exchange policies.
\end{restatable}  

% Shortly, we will see the benefit of working with Definition~\ref{def:exchange-policy}
% instead of Definition~\ref{def:exchange-alg}.

The 
following result establishes an equivalence between exchange algorithms (Definition~\ref{def:exchange-alg}) and exchange policies (Definition~\ref{def:exchange-policy}).
As we will see, this will enable us to leverage tools from optimal stopping to design a welfare-optimal exchange algorithm.

% , shows that we 
% can move back and forth between talking about exchange algorithms and exchange policies.
% Concretely, given an exchange algorithm we can find an equivalent exchange policy and 
% vice versa.
% given an exchange policy we can find an equivalent exchange algorithm.

\begin{restatable}{theorem}{equiv-alg-defs}
    \label{thm:equiv-alg-defs}
    Fix an exchange algorithm
    \(
        \rbr{
            \cbr{
                \algstatemapt
            }_{\timeind\geq 0}, 
            \cbr{
                \algstopmapt
            }_{\timeind\geq 0}
        }
        \in
        \classalgs
    \).
    % as defined in Definition~\ref{def:exchange-alg}.
    Then, we can construct an exchange policy
    \(
        \rbr{
            \cbr{
                \strategyt
            }_{\timeind\geq 0},
            \stoptime 
        }
        \in
        \setexchangepolicies
    \) 
    such that 
    % satisfying Definition~\ref{def:exchange-alg-as-policy} such that
    % \begin{align*}
    \(
        \EE\sbr{
            \rewardss{
                \algstatett{\algendtime}
            }
        }
        =
        \EE\sbr{
            \rewardss{
                \strategytt{\stoptime}
            }
        }
    \).
    % \end{align*}
    Conversely, if given an exchange policy
    % as defined in Definition~\ref{def:exchange-alg-as-policy}.
    \(
        \rbr{
            \cbr{
                \strategyt
            }_{\timeind\geq 0},
            \stoptime 
        }
        \in
        \setexchangepolicies
    \) 
    we can construct an exchange algorithm
    \(
        \rbr{
            \cbr{
                \algstatemapt
            }_{\timeind\geq 0}, 
            \cbr{
                \algstopmapt
            }_{\timeind\geq 0}
        }
        \in
        \classalgs
    \)
    % satisfying Definition~\ref{def:exchange-alg} 
    such that 
    \(
        \EE\sbr{
            \rewardss{
                \strategytt{\stoptime}
            }
        }
        =
        \EE\sbr{
            \rewardss{
                \algstatett{\algendtime}
            }
        }
    \).
    % In other words, exchange algorithms are equivalent to exchange policies.
\end{restatable}  

\begin{proof}
    We start by constructing an exchange policy  
    given an algorithm
    \(
        \rbr{
            \cbr{
                \algstatemapt
            }_{\timeind\geq 0}, 
            \cbr{
                \algstopmapt
            }_{\timeind\geq 0}
        }
    \).
    Define 
    % the sequence
    \(
        \rbr{
            \cbr{
                \strategyt
            }_{\timeind\geq 0},
            \stoptime
        }
    \) 
    via
    \(
        \strategytt{0} = \zerotuple
    \)
    and
    \begin{align*}
    % \(
        \qquad
        \strategytt{\timeind+1}
        =
        \begin{cases}
            \algstatemapt\rbr{
                \cbr{
                    \valaii{k}
                }_{k=0}^{\numdataforatt{t}}
                ,
                \cbr{
                    \valbjj{k}
                }_{k=0}^{\numdataforbtt{t}}
            }
            & 
            \text{ if } \timeind < \algendtime
            \\
            \rbr{
                \numdataa,\numdatab
            }
            &
            \text{ if } \timeind \geq \algendtime
        \end{cases}  
    \end{align*}  
    along with setting
    \(
    % \begin{align*}
        \stoptime
        =
        \algendtime
        % =
        % \inf\cbr{
        %     t\in\NN_{\geq 0} :
        %     \delta_{\timeind}\rbr{
        %         \cbr{
        %             \valaii{k}
        %         }_{k=0}^{\numdataforatt{t}}
        %         ,
        %         \cbr{
        %             \valbjj{k}
        %         }_{k=0}^{\numdataforbtt{t}}
        %     }=0
        % }
    \).

    % \end{align*}  
    We first argue
    \(
        \rbr{
            \cbr{
                \strategyt
            }_{\timeind\geq 0},
            \stoptime
        }
    \) 
    is a well defined policy for 
    \(
        \osasstates=\setexchanges
    \)
    with the family of $\sigma$-algebras 
    \(
        \cbr{
            \statefilts
        }_{\exchange\in\setexchanges}
    \)
    given in~\eqref{eq:family-sig-alg}.
    After we show 
    \(
        \rbr{
            \cbr{
                \strategyt
            }_{\timeind\geq 0},
            \stoptime
        }
    \) 
    satisfies our additional two requirements to be an exchange policy

    \begin{enumerate}
        
    \item
        ($\strategyt\in\setexchanges$): 
        By definition
        \(
            \strategytt{0}=\zerotuple\in\setexchanges
        \).
        For $\timeind\geq 0$, we want to show
        \(
            \strategytt{\timeind+1}\in\setexchanges
        \).
        If $\timeind\geq \algendtime$ then automatically
        \(
            \strategytt{\timeind+1}
            =
            \rbr{
                \numdataa,\numdatab
            }\in\setexchanges
        \).
        If $\timeind < \algendtime$, then $\timeind+1\leq \algendtime$ and
        \[
            \strategytt{\timeind+1}
            =
            \algstatemapt\rbr{
                \cbr{
                    \valaii{k}
                }_{k=0}^{\numdataforatt{t}}
                ,
                \cbr{
                    \valbjj{k}
                }_{k=0}^{\numdataforbtt{t}}
            }
            =
            \algstatett{\timeind+1} .
            % \rbr{
            %     \numdataforatt{\timeind+1},
            %     \numdataforbtt{\timeind+1}
            % } .
        \]
        From Definition~\ref{def:exchange-alg}, we have
        \(
            \algstatett{k}
            <
            \algstatett{k+1}
            % \rbr{
            %     \numdataforatt{k},
            %     \numdataforbtt{k}
            % }
            % <
            % \rbr{
            %     \numdataforatt{k+1},
            %     \numdataforbtt{k+1}
            % }
        \)
        and 
        \(
            % \rbr{
            %     \numdataforatt{\algendtime},
            %     \numdataforbtt{\algendtime}
            % }
            \algstatett{\algendtime}
            \leq
            \rbr{
                \numdataa,\numdatab
            }
        \). 
        Therefore, 
        \(
            \strategytt{\timeind+1}
            =
            \algstatett{\timeind+1}
            \leq
            \algstatett{\algendtime}
            \leq
            \rbr{
                \numdataa,\numdatab
            }
        \)
        and so 
        \(
            \strategytt{\timeind+1}
            \in\setexchanges
        \).

        \item
        ($\strategyt\leq\strategytt{\timeind+1}$):
        For $\timeind\geq \algendtime$, $\strategytt{\timeind+1}$ is constant so this
        holds automatically. For $\timeind<\algendtime$,
        Definition~\ref{def:exchange-alg} requires
        \(
            \strategyt
            =
            \algstatet
            % \rbr{
            %     \numdataforat,
            %     \numdataforbt
            % }
            <
            % \rbr{
            %     \numdataforatt{\timeind+1},
            %     \numdataforbtt{\timeind+1}
            % }
            \algstatett{\timeind+1}
            =
            \strategytt{\timeind+1}
        \). 
    
        \item 
        ($\strategytt{\timeind+1}$
        is
        \(
            \statefiltss{
                \strategyt
            }
        \)-measurable): 
        Notice that we can rewrite $\strategytt{\timeind+1}$ as
        \[
            \strategytt{\timeind+1}
            =
                1_{
                    \cbr{
                        \timeind < \algendtime
                    }
                }
                % \cdot
                \algstatemapt\rbr{
                    \cbr{
                        \valaii{k}
                    }_{k=0}^{\numdataforatt{t}}
                    ,
                    \cbr{
                        \valbjj{k}
                    }_{k=0}^{\numdataforbtt{t}}
                }
                +
                1_{
                    \cbr{
                        \timeind \geq \algendtime
                    }
                }
                % \cdot
                \rbr{
                    \numdataa,\numdatab
                } .
        \]
        We claim that 
        \(
            \cbr{\timeind<\algendtime}
            \in
            \statefiltss{
                \strategyt
            }
        \). To see this, observe that
        \[
            \cbr{\timeind<\algendtime}
            =
            \bigcap_{\ell=0}^\timeind
            \cbr{
                \algstopmaptt{\ell}\rbr{
                    \cbr{
                        \valaii{k}
                    }_{k=0}^{\numdataforatt{\ell}}
                    ,
                    \cbr{
                        \valbjj{k}
                    }_{k=0}^{\numdataforbtt{\ell}}
                }
                = 1
            } .
        \]
        and from the definition of the filtration in
        ~\eqref{eq:family-sig-alg}
        \[
            \sigma\rbr{
                \cbr{
                    \valaii{k}
                }_{k=0}^{\numdataforatt{\ell}}
                ,
                \cbr{
                    \valbjj{k}
                }_{k=0}^{\numdataforbtt{\ell}}
            }
            =
            \statefiltss{
                \rbr{
                    \numdataforatt{\ell},
                    \numdataforbtt{\ell}
                }
            }
        \]
        so each respective event in the union is 
        \(
            \statefiltss{
                \rbr{
                    \numdataforatt{\ell},
                    \numdataforbtt{\ell}
                }
            }
        \)-measurable. Since $\algstatett{1}\leq\algstatett{2}$ if and only if 
        \(
            \statefiltss{
                \algstatett{1}
            }
            \subseteq
            \statefiltss{
                \algstatett{2}
            }
        \), we have 
        \(
            \statefiltss{
                \rbr{
                    \numdataforatt{\ell},
                    \numdataforbtt{\ell}
                }
            }
            \subseteq
            \statefiltss{
                \rbr{
                    \numdataforat,
                    \numdataforbt
                }
            }
            =
            \statefiltss{
                \strategyt
            }
        \), which implies
        \(
            \cbr{
                % \stoptime\leq\timeind
                \timeind < \algendtime
            }
        \)
        is 
        \(
            \statefiltss{
                \strategyt
            }
        \)-measurable.
        Since 
        \[
            \sigma\rbr{
                \cbr{
                    \valaii{k}
                }_{k=0}^{\numdataforat}
                ,
                \cbr{
                    \valbjj{k}
                }_{k=0}^{\numdataforbt}
            }
            =
            \statefiltss{
                \rbr{
                    \numdataforat,
                    \numdataforbt
                }
            },
        \]
        we get 
        \(
            \algstatemapt\rbr{
                \cbr{
                    \valaii{k}
                }_{k=0}^{\numdataforatt{t}}
                ,
                \cbr{
                    \valbjj{k}
                }_{k=0}^{\numdataforbtt{t}}
            }
        \)
        is 
        \(
            \statefiltss{
                \rbr{
                    \numdataforat,
                    \numdataforbt
                }
            }
        \)-measurable, thus 
        $\strategytt{\timeind+1}$ is too as the composition and product of measurable functions.


        % \mycomment{need to argue that the selection is also measurable}

        % This holds automatically for $\timeind\geq\algendtime$ since 
        % $\strategytt{\timeind+1}$ is constant. For $\timeind<\algendtime$,
        % $\strategytt{\timeind+1}$ is a measurable function of 
        % \(
        %     \rbr{
        %         \cbr{
        %             \valaii{k}
        %         }_{k=0}^{\numdataforatt{t}}
        %         ,
        %         \cbr{
        %             \valbjj{k}
        %         }_{k=0}^{\numdataforbtt{t}}
        %     }
        % \).
        % This means $\strategytt{\timeind+1}$ is 
        % \(
        %     \sigma\rbr{
        %         \cbr{
        %             \valaii{k}
        %         }_{k=0}^{\numdataforatt{t}}
        %         ,
        %         \cbr{
        %             \valbjj{k}
        %         }_{k=0}^{\numdataforbtt{t}}
        %     }
        % \)-measurable.
        % Moreover, from the definition of the filtration in~\eqref{eq:family-sig-alg},
        % \[
        %     \sigma\rbr{
        %         \cbr{
        %             \valaii{k}
        %         }_{k=0}^{\numdataforatt{t}}
        %         ,
        %         \cbr{
        %             \valbjj{k}
        %         }_{k=0}^{\numdataforbtt{t}}
        %     }
        %     =
        %     \statefiltss{
        %         \rbr{
        %             \numdataforat,\numdataforbt
        %         }
        %     }   
        %     =
        %     \statefiltss{
        %         \strategyt
        %     } .
        % \]
        % % so $\strategytt{\timeind+1}$
        % % is
        % % \(
        % %     \statefiltss{
        % %         \strategyt
        % %     }
        % % \)-measurable

        \item
        \big($\stoptime$ is a stopping time w.r.t. the filtration
        \(
            \cbr{
                \timefiltt
            }_{\timeind\geq 0}
            =
            \cbr{
                \statefiltss{
                    \strategyt
                }
            }_{\timeind\geq 0}
        \)\big):
        Start by rewriting
        \[
            \cbr{
                \stoptime\leq\timeind
            }
            =
            \bigcup_{\ell=0}^\timeind
            \cbr{
                \algstopmaptt{\ell}\rbr{
                    \cbr{
                        \valaii{k}
                    }_{k=0}^{\numdataforatt{\ell}}
                    ,
                    \cbr{
                        \valbjj{k}
                    }_{k=0}^{\numdataforbtt{\ell}}
                }
                = 0
            } .
        \]
        Now notice that 
        \[
            \sigma\rbr{
                \cbr{
                    \valaii{k}
                }_{k=0}^{\numdataforatt{\ell}}
                ,
                \cbr{
                    \valbjj{k}
                }_{k=0}^{\numdataforbtt{\ell}}
            }
            =
            \statefiltss{
                \rbr{
                    \numdataforatt{\ell},
                    \numdataforbtt{\ell}
                }
            }
        \]
        so each respective event in the union is 
        \(
            \statefiltss{
                \rbr{
                    \numdataforatt{\ell},
                    \numdataforbtt{\ell}
                }
            }
        \)-measurable. Since $\algstatett{1}\leq\algstatett{2}$ if and only if 
        \(
            \statefiltss{
                \algstatett{1}
            }
            \subseteq
            \statefiltss{
                \algstatett{2}
            }
        \), we have 
        \(
            \statefiltss{
                \rbr{
                    \numdataforatt{\ell},
                    \numdataforbtt{\ell}
                }
            }
            \subseteq
            \statefiltss{
                \rbr{
                    \numdataforat,
                    \numdataforbt
                }
            }
            =
            \statefiltss{
                \strategyt
            }
        \), which implies
        \(
            \cbr{
                \stoptime\leq\timeind
            }
        \)
        is 
        \(
            \statefiltss{
                \strategyt
            }
            =
            \timefiltt
        \)-measurable.
        Therefore, 
        $\stoptime$ is a stopping time with respect to the filtration
        \(
            \cbr{
                \timefiltt
            }_{\timeind\geq 0}
        \).
        
    \end{enumerate}  

    We are left with the two additional requirements of 
    Definition~\ref{def:exchange-policy}. For the first requirement, suppose 
    that
    \(
        \strategyt <
        \rbr{
            \numdataa, \numdatab
        }
    \), then by construction this implies $\timeind<\algendtime=\stoptime$
    so $\strategyt=\algstatet<\algstatett{\timeind+1}=\strategytt{\timeind+1}$
    by Definition~\ref{def:exchange-alg}.
    For the second requirement, we know that $\stoptime=\algendtime$ is finite or else 
    it would contradict the fact that 
    \(
        \algstatett{\algendtime}
        \leq
        \rbr{
            \numdataa,\numdatab
        }
    \), since $\cbr{\algstatet}_{\timeind\geq 0}$ is increasing.
    From our construction, we have that 
    \(
        \rbr{
            \algstatett{0},
            \ldots,
            \algstatett{\algendtime}
        }
        =
        \rbr{
            \strategytt{0},
            \ldots,
            \strategytt{\stoptime}
        }
    \), therefore
    \(
        \EE\sbr{
            \rewardss{
                \strategytt{\stoptime}
            }
        }
        =
        \EE\sbr{
            \rewardss{
                \algstatett{\algendtime}
            }
        }
    \).

    For the reverse direction, we create an exchange algorithm given an
    exchange policy 
    \(
        \rbr{
            \cbr{
                \strategyt
            }_{\timeind\geq 0},
            \stoptime 
        }
    \).
    Let 
    \(
        \rbr{
            \numdataforat,
            \numdataforbt
        }
        =
        \strategyt
    \).
    From Definition~\ref{def:exchange-policy} we know $\stoptime$ is a stopping time,
    thus
    \(
        1_{
            \cbr{
                \stoptime\neq\timeind
            }
        }
    \)
    is
    \[
        \statefiltss{\strategyt}
        =
        \sigma\rbr{
            \cbr{
                \valaii{k}
            }_{k=0}^{\numdataforat}
            ,
            \cbr{
                \valbjj{k}
            }_{k=0}^{\numdataforbt}
        }
    \]measurable.
    % is 
    % \(
    % \)-measurable, 
    Therefore,
    there exists a measurable function $\algstopmapt$ such that 
    \[
        1_{
            \cbr{
                \stoptime\neq\timeind
            }
        }
        =
        \algstopmapt\rbr{
            \cbr{
                \valaii{k}
            }_{k=0}^{\numdataforat}
            ,
            \cbr{
                \valbjj{k}
            }_{k=0}^{\numdataforbt}
        }.
    \]
    Similarly, since $\strategytt{\timeind+1}$ is 
    \(
        \statefiltss{\strategyt}
    \)-measurable,
    there exists a measurable function 
    \(
        \mbleexistsalgstatemapt
    \)
    such that
    \[
        \strategytt{\timeind+1}
        =
        \mbleexistsalgstatemapt
        \rbr{
            \cbr{
                \valaii{k}
            }_{k=0}^{\numdataforat}
            ,
            \cbr{
                \valbjj{k}
            }_{k=0}^{\numdataforbt}
        } .
    \]
    Now define 
    \(
        \cbr{
            \algstatemapt
        }_{\timeind\geq 0}
    \) 
    via
    \[
        \algstatemapt
        =
        \begin{cases}
            \mbleexistsalgstatemapt
            &
            \text{if }~~~
            \algstatet < 
            \rbr{
                \numdataa,\numdatab
            }
            % \forall r\in\cbr{0,\ldots,\timeind}
            % \algstopmapt\rbr{
            %     \cbr{
            %         \valaii{k}
            %     }_{k=0}^{\numdataforatt{r}}
            %     ,
            %     \cbr{
            %         \valbjj{k}
            %     }_{k=0}^{\numdataforbtt{r}}
            % }= 0
            \\
            \algstatet+(1,1)
            & 
            \text{otherwise}
        \end{cases} .
    \]
    By the requirement that
    \(
        \strategyt < \strategytt{\timeind+1}
    \)
    for 
    \(
        \strategyt <
        \rbr{
            \numdataa,\numdatab
        }
    \), we get
    $\algstatet <\algstatett{\timeind+1}$.
    Moreover, by our construction
    \(
        \rbr{
            \algstatett{0},
            \ldots,
            \algstatett{\algendtime}
        }
        =
        \rbr{
            \strategytt{0},
            \ldots,
            \strategytt{\stoptime}
        },
    \)
    % we also have that 
    which implies
    \(
        \algstatett{\algendtime}
        \leq
        \rbr{
            \numdataa,\numdatab
        }
    \)
    and 
    \(
        \EE\sbr{
            \rewardss{
                \strategytt{\stoptime}
            }
        }
        =
        \EE\sbr{
            \rewardss{
                \algstatett{\algendtime}
            }
        }
    \).
\end{proof}  

Therefore, going forward, we will refer to exchange algorithms and exchange policies interchangeably.

% \mycomment{delete below ->}
% \begin{restatable}{definition}{exchange-alg-as-policy2}
%     \label{def:exchange-alg-as-policy2} 
%     A resource exchange algorithm is a tuple 
%     \(
%         \rbr{
%             \cbr{
%                 \strategyt
%             }_{\timeind\geq 0},
%             \stoptime 
%         }
%     \).
%     Here
%     \(
%         \strategyt\in\setexchanges
%     \)
%     is a non-decreasing sequence
%     satisfying
%     \begin{align*}
%         \strategytt{0}=0,
%         \quad
%         \strategytt{\timeind+1}
%         \text{ is }
%         \statefiltss{
%             \strategyt
%         }\text{-measurable},
%         \quad
%         % \text{and }
%         \strategyt < \strategytt{\timeind+1}
%         \text{ for }
%         \strategyt < \rbr{
%             \numdataa,\numdatab
%         },
%     \end{align*}  
%     and $\stoptime$ is a finite stopping time with respect to the filtration 
%     \(
%         \cbr{
%             \timefiltt
%         }_{\timeind\geq 0}
%         =
%         \cbr{
%             \statefilt_{
%                 \strategyt
%             }
%         }_{\timeind\geq 0}
%     \)
%     as given in~\eqref{eq:family-sig-alg}.
% \end{restatable}  
% \mycomment{delete above <-}

% Definition~\ref{def:exchange-alg-as-policy} is exactly the definition of a policy in~\citep{mandelbaum1981optimal}
% where the state space is $\setexchanges$, and the family of $\sigma$-algebras over the state space is 
% \(
%     \cbr{
%         \statefilts
%     }_{\exchange\in\setexchanges}
% \).
% If, in the language of~\citep{mandelbaum1981optimal}, 

\subsection{Reducing Optimal Exchange Algorithm Design to Optimal Stopping}
\label{subsec:reduction-to-optimal-stopping}

We will now show that 
by defining the reward for an exchange $\exchange\in\setexchanges$ to be 
$\rewards$, 
the problem of designing an optimal exchange algorithm $\alg\in\classalgs$,
will reduce to solving an optimal stopping time problem over $\setexchanges$.
Before doing this, we first introduce a 
% few necessary concepts 
necessary concept
from optimal stopping theory.
% and~\citet{mandelbaum1981optimal}.

\begin{restatable}{definition}{partially-ordered-snell-env}
    \textup{
        (Snell envelope, \citet{mandelbaum1981optimal})
    }
    For any $\osasstate\in\osasstates$, let
    \[
        \osassnellenvs
        =
        \esssup_{
            \osaspolicy
            \in
            \osaspoliciess
        }
        \EE\sbr{
            \osasrewardss{
                \osasstrategytt{
                    \osasstoptime
                }
            }
            | \osasstatefilts
        }
    \]
    where 
    $\osaspoliciess$ is the set of policies for which 
    \(
        \osasstrategytt{0} = \osasstate
    \).
    The Snell envelope is defined as the family 
    \(
        \cbr{
            \osassnellenvs
        }_{
            \osasstate\in\osasstates
        }
    \).
\end{restatable}  

Intuitively, each $\osassnellenvs$ represents the best expected reward, given the information available
at $\osasstate$, following policies that start at $\osasstate$. The definition of the Snell envelope 
depends on the type of optimal stopping problem considered; however, it is usually difficult to calculate 
$\osassnellenvs$ directly from the definiton. To remidy this, a dymanic programming equation is often
proved. 
This equation usually takes the following form: $\osassnellenvs$ is the maximum of stopping and accepting
the current reward or the expected reward of ``proceeding optimally''.
Using this, we now present an optimal exchange algorithm.


% The idea is to show that $\osassnellenvs$ is the maximum of the following two choices: stop and accept the 
% current reward or 

% The following is such a result in the case of optimal stopping times over partially ordered sets.

% \mycomment{move this to appendix but cite here, describe high level}
% \kkcomment{do this.}

% \begin{restatable}{theorem}{osas-dp-eq}
%     \label{thm:osas-dp-eq}
%     \textup{
%         (\citet{mandelbaum1981optimal})
%     }
%     Let 
%     \(
%         \osasdirectsuccessorss{
%             \osasstate
%         }
%         =
%         \cbr{
%             u> \osasstate: \sbr{\osasstate,u}=\cbr{\osasstate,u}
%         } 
%     \)
%     be the set of direct successors for $\osasstate$, then
%     \begin{align*}
%         \osassnellenvs
%         =
%         \max
%         \cbr{
%             \osasrewards,
%             \EE\sbr{
%                 \osassnellenvss{u}
%                 | 
%                 \osasstatefilts
%             }
%             : u\in
%             \osasdirectsuccessorss{
%                 \osasstate
%             }
%         }
%         \quad
%         \text{and}
%         \quad
%         \EE\sbr{
%             \osassnellenvss{0}
%         }
%         =
%         \sup_{
%             \osaspolicy\in\osaspolicies
%         }
%         \EE\sbr{
%             \osasrewardss{
%                 \osasstrategytt{
%                     \osasstoptime
%                 }
%             }
%         }
%         .
%     \end{align*}  
% \end{restatable}

% The first part of this result says that the best expected reward of a policy which starts from state
% $\osasstate\in\osasstates$ is the maximum of the reward available from stopping at $\osasstate$,
% $\osasrewards$, and the best 
% expected reward among policies starting at a successor $u$, conditioned on the information 
% available at $\osasstate$. The second part shows that $\osassnellenvss{0}$ characterizes the 
% performance of an optimal policy. Using this, we now present an optimal exchange algorithm,

% \mycomment{remove all the tildes? Unnecessary clutter? I did this with the IC section}

% \kkcomment{Since this is a key result of the paper, you can also move the IR rsult and proof into this theorem. Currently, it looks very disjointed. Also, better to state this in a way where you are presenting the }

\begin{restatable}{theorem}{opt-alg}
    \label{thm:opt-alg}
    % In the setting of~\citep{mandelbaum1981optimal}, 
    Let $\osasstates=\setexchanges$, $\osasrewards=\rewards$,
    and the family 
    \(
        \cbr{
            \statefilts
        }_{\exchange\in\setexchanges}
    \) be given by~\eqref{eq:family-sig-alg}.
    Define
    the policy
    \(
        \optexpolicy=
        \rbr{
            \cbr{
                \optstrategyt
            }_{\timeind\geq 0},
            \optstoptime
        }
    \)
    via
    %  $\strategytt{0}=0$, for $\timeind\geq 1$
    \begin{align*}
        \optstrategytt{0} = (0,0),
        \quad
        % \text{for } \timeind\geq 1
        \optstrategytt{\timeind+1}
        =
        \argmax_{
            \nextexchange\in 
            \directsuccessorss{
                \optstrategyt
            }
        }
        \EE\sbr{
            \snellenvss{\nextexchange}
            | \timefiltt
        }
        ~~~\text{ for }~~~~\timeind\geq 0,
        % \quad
        % \text{and}
        % \quad
        % \stoptime
    \end{align*}  
    % where 
    % \(
    %     \directsuccessorss{
    %         % \optstrategyt
    %         \exchange
    %     }
    %     =
    %     \cbr{
    %         \exchange + (1,0),
    %         \exchange + (0,1)
    %     }
    % \),
    and 
    \(
        \optstoptime
        =
        \inf
        \cbr{
            t\in\NN_{\geq 0} :
            \snellenvss{
                \optstrategyt
            }
            =
            \rewardss{
                \optstrategyt
            }
        }
    \). 
    Then $\optexpolicy\in\setexchangepolicies$ and is optimal, i.e.
    % \begin{align*}
    \(
        \EE\sbr{
            \rewardss{
                \optstrategytt{\optstoptime}
            }
        }
        =
        \sup_{
            % \rbr{
            %     \cbr{\strategyt}_{\timeind\geq 0},
            %     \stoptime
            % }
            % \in\classalgs
            \exchangepolicy
            \in
            \setexchangepolicies
        }
        \EE\sbr{
            \rewardss{
                \strategytt{\stoptime}
            }
        }
    \). 
    Moreover if
    for the choice 
    of welfare function $\welfarefunc$,
    \[
        \EE\sbr{
            \welfarefunc\rbr{
                \utilass{
                    \optstrategytt{\optstoptime}
                },
                \utilbss{
                    \optstrategytt{\optstoptime}
                }
            }
        }
        \geq 0
        ~~~
        \Rightarrow
        ~~~
        \EE\sbr{
            \utilass{
                \optstrategytt{\optstoptime}
            }
        },
        \EE\sbr{
            \utilbss{
                \optstrategytt{\optstoptime}
            }
        }
        \geq 0
    \]
    then $\optexpolicy$ is individually rational.
\end{restatable}  
\begin{proof}
    We first verify that $\optexpolicy$ is a well defined policy.
    Since 
    \(
        \timefiltt
        =
        \statefiltss{
            \optstrategyt
        }
    \),
    \(
        \optstrategytt{\timeind+1}
    \)
    is the argmax of a finite number of
    \(
        \statefiltss{
            \optstrategyt
        }
    \)-measurable random variables, 
    and thus 
    \(
        \optstrategytt{\timeind+1}
    \)
    is
    \(
        \statefiltss{
            \optstrategyt
        }
    \)-measurable.
    From the definition of 
    \(
        \directsuccessorss{
            \optstrategyt
        }
    \), it follows that
    \(
        \optstrategyt
        \leq
        \optstrategytt{\timeind+1}
    \).
    Now observe that 
    % \begin{align*}
    \(
        \cbr{
            \optstoptime\leq \timeind
        }
        =
        \bigcup_{k=0}^\timeind
        \cbr{
            \snellenvss{
                \optstrategytt{k}
            }
            =
            \rewardss{
                \optstrategytt{k}
            }
        }
    \)
    and 
    \(
        \cbr{
            \snellenvss{
                \optstrategytt{k}
            }
            =
            \rewardss{
                \optstrategytt{k}
            }
        }
        \in
        \timefilttt{k}
        \subseteq
        \timefiltt
    \)
    by the definition in~\eqref{eq:family-sig-alg}.
    Thus
    \(
        \cbr{
            \optstoptime\leq \timeind
        }
        \in
        \timefiltt
    \)
    so $\optstoptime$ is a well defined stopping time with respect to 
    \(
        \cbr{
            \timefiltt
        }_{\timeind\geq 0}
    \).

    We now argue $\optexpolicy$ is a valid exchange policy. Since at time $\timeind\geq 0$ the next exchange
    is selected from 
    \(
        \directsuccessorss{
            \optstrategyt
        }
    \),
    we have
    \(
        \optstrategyt < \optstrategytt{\timeind+1} 
    \).
    Now suppose that $\optstoptime$ is not finite. Then, for sufficiently a sufficiently large $T$,
    we would have 
    \(
        \optstrategytt{T}
        =
        \rbr{
            \numdataa,\numdatab
        }
    \).
    Since 
    \(
        \directsuccessorss{
            \optstrategytt{T}
        }
        =
        \varnothing
    \),
    Theorem~\ref{thm:osas-dp-eq}
    implies that 
    \(
        \snellenvss{
            \optstrategytt{T}
        }
        =
        \rewardss{
            \optstrategytt{T}
        }
    \), so $\optstoptime$ is forced to be finite.

    By construction, $\optexpolicy$
    is the natural policy which satisfies conditions~\eqref{eq:osas-thm2-cond1}--\eqref{eq:osas-thm2-cond3}.
    %  (2.5)--(2.7) in~\citep{mandelbaum1981optimal} 
    %  \kkcomment{State the theorems you are using in the appendix and refer to the conditions here.}. 
    Therefore, together with the fact that $\optstoptime$ is finite, we can apply 
    Theorem~\ref{thm:osas-thm2}
    % Theorem 2
    % from~\citep{mandelbaum1981optimal} 
    to conlude $\optexpolicy$ is an optimal policy, i.e.
    \(
        \EE\sbr{
            \rewardss{
                \optstrategytt{\optstoptime}
            }
        }
        =
        \sup_{
            \apolicy\in
            \osaspolicies
        }
        \EE\sbr{
            \rewardss{
                \strategytt{\stoptime}
            }
        }
    \).
    Since $\setexchangepolicies\subseteq\osaspolicies$, and $\optexpolicy\in\setexchangepolicies$ 
    \[
        \sup_{
            \apolicy\in
            % \classalgs
            \setexchangepolicies
        }
        \EE\sbr{
            \rewardss{
                \strategytt{\stoptime}
            }
        }
        \leq
        \sup_{
            \apolicy\in
            \osaspolicies
        }
        \EE\sbr{
            \rewardss{
                \strategytt{\stoptime}
            }
        }
        =
        \EE\sbr{
            \rewardss{
                \optstrategytt{\optstoptime}
            }
        }
        \leq
        \sup_{
            \apolicy\in
            \setexchangepolicies
        }
        \EE\sbr{
            \rewardss{
                \strategytt{\stoptime}
            }
        } .
    \]
    Therefore
    \(
        \EE\sbr{
            \rewardss{
                \optstrategytt{\optstoptime}
            }
        }
        =
        \sup_{
            % \rbr{
            %     \cbr{\strategyt}_{\timeind\geq 0},
            %     \stoptime
            % }
            % \in\classalgs
            \exchangepolicy
            \in
            \setexchangepolicies
        }
        \EE\sbr{
            \rewardss{
                \strategytt{\stoptime}
            }
        }
    \),
    so $\optexpolicy$
    is an optimal exchange policy.

    For the second part of the claim,
    consider the exchange policy which facilitates no exchanges, i.e.
    \(
        \altexchangepolicy
        =
        \rbr{
            \altstrategytt{0} = (0,0),
            \altstoptime = 0
        } 
    \).
    Then by our normalization assumptions we have
    \(
        \valaii{0} = \valbjj{0} = 0
    \)
    which implies 
    \(
        \utilass{0}
        =
        \utilbss{0}
        =
        \welfarefunc\rbr{
            \utilass{0},
            \utilbss{0}
        }
        =0
    \).
    Therefore, since $\optexpolicy$ is optimal,
    \(
        0
        =
        \EE\sbr{
            \rewardss{
                \altstrategytt{\altstoptime}
            }
        }
        \leq
        \EE\sbr{
            \rewardss{
                \optstrategytt{\optstoptime}
            }
        }
        =
        \EE\sbr{
            \welfarefunc\rbr{
                \utilass{
                    \optstrategytt{\optstoptime}
                },
                \utilbss{
                    \optstrategytt{\optstoptime}
                }
            }
        }
    \),
    so our assumption yields individual rationality.
\end{proof}  

It is worth pointing out that there are reasonable choices of welfare functions that satisfy the implication
in the Theorem~\ref{thm:opt-alg}. 
A simple example of such a welfare function is the egalitarian welfare
\(
    \welfarefunc\rbr{x,y}=\min(x,y)
\)
since for any random variables $X$ and $Y$,
\(
    \EE\sbr{
        \min(X,Y)
    }
    \geq 0
    \Rightarrow
    \EE\sbr{X},
    \EE\sbr{Y}
    \geq 0
\).


\subsection{An example}
\label{sec:algoexample}

We now revisit Example~\ref{ex:multi-vs-single-round-alg} to demonstrate how Theorem~\ref{thm:opt-alg}
is gives an optimal algorithm.
We will show that the optimal exchange algorithm is given by
    \(
        \rbr{
            \cbr{\algstatemapt}_{\timeind=0}^{1},
            \cbr{\algstopmap}_{\timeind=0}^{2}
        }
    \)
    where 
    \begin{gather*}
        \algstatemaptt{0}(0,0) = (1,0)
        \qquad
        \algstatemaptt{1}(1,0) = (1,1)
        \\
        \algstopmaptt{0}(0,0) = 1
        \qquad
        \algstopmaptt{1}(1,0) 
        =
        \begin{cases}
            1 & \text{if } ~\valaii{1} = 1
            \\
            0 & \text{if } ~\valaii{1} = 0
        \end{cases}  
        \qquad
        \algstopmaptt{2}(1,1) 
        = 0 . 
    \end{gather*}  
    and moreover that,
    \(
        \EE\sbr{
            \rewardss{
                \algstatett{\algendtime}
            }
        }
        =
        \frac{1}{8}
    \).

% \begin{restatable}{claim}{two-bern-ex-opt-alg}
%     \label{ex:two-bern-ex-opt-alg}
%     % Let $\numdataa=\numdatab=1$.
%     % Suppose that each competitor's value for the other's good is 
%     % $\valaii{1},\valbjj{1}\drawniid\Bern(1/2)$ and assume $\valaii{0}=\valbjj{0}=0$.
%     % Let the utilities and welfare be given by
%     % \begin{align*}
%     %     \utilas
%     %     =
%     %     \valai-\frac{1}{4}\valbj,
%     %     \quad 
%     %     \utilbs
%     %     =
%     %     \valbj-\frac{1}{4}\valai,
%     %     \quad
%     %     \text{and}
%     %     \quad
%     %     \welfarefunc(x,y)=\min(x,y) .
%     % \end{align*}  
%     Recall the environment from Example~\ref{ex:two-bern-ex-opt-alg}.
%     The optimal exchange algorithm is given by
%     \(
%         \rbr{
%             \cbr{\algstatemapt}_{\timeind=0}^{1},
%             \cbr{\algstopmap}_{\timeind=0}^{2}
%         }
%     \)
%     where 
%     \begin{gather*}
%         \algstatemaptt{0}(0,0) = (1,0)
%         \qquad
%         \algstatemaptt{1}(1,0) = (1,1)
%         \\
%         \algstopmaptt{0}(0,0) = 0
%         \qquad
%         \algstopmaptt{1}(1,0) 
%         =
%         \begin{cases}
%             1 & \text{if } ~\valaii{1} = 0
%             \\
%             0 & \text{if } ~\valaii{1} = 1
%         \end{cases}  
%         \qquad
%         \algstopmaptt{2}(1,1) 
%         = 1 . 
%     \end{gather*}  
%     and 
%     \(
%         \EE\sbr{
%             \rewardss{
%                 \algstatett{\algendtime}
%             }
%         }
%         =
%         \frac{1}{8}
%     \).
% \end{restatable}  

\begin{proof}
    We will use Theorem~\ref{thm:opt-alg} to determine the optimal exchange policy which will correspond to 
    the optimal exchange algorithm above. 
    % \mycomment{add and thus is the opt alg by the corollary to add?}
    To determine the optimal exchange policy, we first calculate the Snell envelope
    \(
        \cbr{
            \snellenvs
        }_{\exchange\in\setexchanges}
    \) working backwards.

    From Theorem~\ref{thm:osas-dp-eq} we have
    % \begin{align*}
    \(
        \snellenvss{(1,1)}
        =
        \rewardss{(1,1)}
    \)
    % \end{align*}  
    which gives
    \begin{align*}
        \snellenvss{(1,0)}
        =
        \max\rbr{
            \rewardss{(1,0)},
            \EE\sbr{
                \snellenvss{(1,1)}
                | \statefiltss{(1,0)}
            }
        }
        =
        \max\rbr{
            \rewardss{(1,0)},
            \EE\sbr{
                \rewardss{(1,1)}
                | \valaii{1}
            }
        }
        \\
        \snellenvss{(0,1)}
        =
        \max\rbr{
            \rewardss{(0,1)},
            \EE\sbr{
                \snellenvss{(1,1)}
                | \statefiltss{(0,1)}
            }
        }
        =
        \max\rbr{
            \rewardss{(0,1)},
            \EE\sbr{
                \rewardss{(1,1)}
                | \valbjj{1}
            }
        }
    \end{align*}  
    and thus by symmetry,
    \begin{align*}
        \snellenvss{(0,0)}
        &=
        \max\rbr{
            \rewardss{(0,0)},
            \EE\sbr{
                \snellenvss{(1,0)}
                | \statefiltss{(0,0)}
            },
            \EE\sbr{
                \snellenvss{(0,1)}
                | \statefiltss{(0,0)}
            }
        }
        \\
        &=
        \max\rbr{
            \rewardss{(0,0)},
            \EE\sbr{
                \snellenvss{(1,0)}
            },
            \EE\sbr{
                \snellenvss{(0,1)}
            }
        }
        \\
        &=
        \max\rbr{
            0,
            \EE\sbr{
                \snellenvss{(1,0)}
            }
        } .
    \end{align*}  
    To determine
    \(
        \EE\sbr{
            \rewardss{(1,1)}
            | \valaii{1}
        }
    \), consider two cases.

        If $\valaii{1}=0$ then $\rewardss{(1,1)}=0$ when $\valbjj{1}=0$, and 
        $\rewardss{(1,1)}=-\frac{1}{4}$ when $\valbjj{1}=1$ so
        \[
            \EE\sbr{
                \rewardss{(1,1)}
                | \valaii{1}=0
            }
            =
            \frac{1}{2}\cdot 0
            +
            \frac{1}{2}\cdot\rbr{-\frac{1}{4}} 
            =
            -\frac{1}{8}
        \]

        If $\valaii{1}=1$ then $\rewardss{(1,1)}=-\frac{1}{4}$ when $\valbjj{1}=0$, and 
        $\rewardss{(1,1)}=\frac{3}{4}$ when $\valbjj{1}=1$ so
        \[
            \EE\sbr{
                \rewardss{(1,1)}
                | \valaii{1}=1
            }
            =
            \frac{1}{2}\cdot \rbr{-\frac{1}{4}}
            +
            \frac{1}{2}\cdot\rbr{\frac{3}{4}} 
            =
            \frac{1}{4} .
        \]

    \noindent
    Therefore, 
    \(
        \EE\sbr{
            \rewardss{(1,1)}
            | \valaii{1}
        }
        =
        \frac{3}{8}\valaii{1}
        -
        \frac{1}{8}
    \).
    We use this to calculate
    \begin{align*}
        \EE\sbr{
            \snellenvss{(1,0)}
        }
        &=
        \max\rbr{
            \rewardss{(1,0)},
            \frac{3}{8}\valaii{1}
            -
            \frac{1}{8}
        }
        \\
        &=
        \frac{1}{4}
        \max\rbr{
            \min\rbr{
                0,0
            },
            -
            \frac{1}{8}
        }
        +
        \frac{1}{4}
        \max\rbr{
            \min\rbr{
                1, -\frac{1}{4}
            },
            \frac{3}{8}-\frac{1}{8}
        }
        \\
        &~+
        \frac{1}{4}
        \max\rbr{
            \min(0,0),
            -\frac{1}{8}
        }
        +
        \frac{1}{4}
        \max\rbr{
            \min\rbr{
                1, -\frac{1}{4}
            },
            \frac{3}{8}-\frac{1}{8}
        }
        \\
        &= 
        \frac{1}{8}
    \end{align*}  
    and thus 
    \(
        \snellenvss{(0,0)}
        =
        \max\rbr{0,\frac{1}{8}}
        =
        \frac{1}{8}
    \).

    By definition, the policy starts at
    \(
        \optstrategytt{0}
        =
        (0,0)
    \).
    Because we define 
    \(
        \optstoptime
        =
        \inf
        \cbr{
            t\in\NN_{\geq 0} :
            \snellenvss{
                \optstrategyt
            }
            =
            \rewardss{
                \optstrategyt
            }
        }
    \),
    we have $\optstoptime>0$, i.e. the policy does not stop immediately and since 
    \[
        \optstrategytt{1}
        =
        \argmax_{
            \nextexchange\in 
            \cbr{
                (1,0),(0,1)
            }
        }
        \EE\sbr{
            \snellenvss{\nextexchange}
        }
    \]
    it proceeds to either $(1,0)$ $(0,1)$. Because of symmetry 
    \(
        \EE\sbr{
            \snellenvss{(1,0)}
        }
        =
        \EE\sbr{
            \snellenvss{(0,1)}
        }
    \)
    so without loss of generality $\optstrategytt{1}=(1,0)$.
    At this point
    \(
        \optstoptime
        =
        1
    \)
    exactly when
    \[
        \rewardss{(1,0)}
        \geq
        \EE\sbr{
            \rewardss{(1,1)}
            | \valaii{1}
        }
        \iff
        \valaii{1}=0,
    \]
    so the policy stops if $\valaii{1}=0$.
    Otherwise, if $\valaii{1}=1$ the policy continues to
    \[
        \optstrategytt{2}
        =
        \argmax_{
            \nextexchange\in 
            \cbr{
                (1,1)
            }
        }
        \EE\sbr{
            \snellenvss{\nextexchange}
        }
        =
        (1,1)
    \]
    as the only remaining state.
    Finally, if the policy continued, 
    then $\optstoptime=2$
    since $\snellenvss{(1,1)}=\rewardss{(1,1)}$.
    Therefore, the $\optexpolicy$ corresponds to the exchange algorithm above.
\end{proof}  

For more complex models it may be difficult to analytically compute the condition expectations
used to determine the sequence of exchanges
\(
    \cbr{
        \optstrategyt
    }_{\timeind\geq 0}
\)
and when to terminate the exchanges $\optstoptime$. This is a common phenomena in optimal 
stopping, in which case numerical approximation methods such as Monte Carlo based techniques are 
used~\citep{longstaff2001valuing, tsitsiklis2001regression, broadie1997pricing}.


% \mycomment{add corollary that we have an opt exchange alg?}

% \subsection{An IR Sufficient Condition}

% \kkcomment{Move up to main theorem.}

% In addition to being optimal, the following result shows that for a simple condition on the welfare function
% $\welfarefunc(x,y)$, the algorithm above is also individually rational.

% \begin{restatable}{proposition}{ir-welfare-condition}
%     Let $\optalg\in\classalgs$ be the algorithm given in Theorem~\ref{thm:opt-alg}. 
%     If for the choice 
%     of welfare function $\welfarefunc$,
%     \[
%         % \EE\sbr{
%         %     \rewardss{
%         %         \optstrategytt{\optstoptime}
%         %     }
%         % }
%         \EE\sbr{
%             \welfarefunc\rbr{
%                 \utilass{
%                     \optstrategytt{\optstoptime}
%                 },
%                 \utilbss{
%                     \optstrategytt{\optstoptime}
%                 }
%             }
%         }
%         \geq 0
%         ~~~
%         \Rightarrow
%         ~~~
%         \EE\sbr{
%             \utilass{
%                 \optstrategytt{\optstoptime}
%             }
%         },
%         \EE\sbr{
%             \utilbss{
%                 \optstrategytt{\optstoptime}
%             }
%         }
%         \geq 0
%     \]
%     then $\optalg$ is individually rational.
% \end{restatable}  
% \begin{proof}
%     Consider the algorithm which facilitates no exchanges, i.e.
%     \(
%         \alg
%         =
%         \rbr{
%             \strategytt{0} = (0,0),
%             \stoptime = 0
%         } 
%     \).
%     Then by our normalization assumptions we have
%     \(
%         \valaii{0} = \valbjj{0} = 0
%     \)
%     which implies 
%     \(
%         \utilass{0}
%         =
%         \utilbss{0}
%         =
%         \welfarefunc\rbr{
%             \utilass{0},
%             \utilbss{0}
%         }
%         =0
%     \).
%     Therefore, since $\optalg$ is optimal,
%     \(
%         0
%         =
%         \EE\sbr{
%             \rewardss{
%                 \strategytt{\stoptime}
%             }
%         }
%         \leq
%         \EE\sbr{
%             \rewardss{
%                 \optstrategytt{\optstoptime}
%             }
%         }
%         =
%         \EE\sbr{
%             \welfarefunc\rbr{
%                 \utilass{
%                     \optstrategytt{\optstoptime}
%                 },
%                 \utilbss{
%                     \optstrategytt{\optstoptime}
%                 }
%             }
%         }
%     \)
%     ,
%     so our assumption yields individual rationality.
% \end{proof}  

% A simple example of such a welfare function is 
% \(
%     \welfarefunc\rbr{x,y}=\min(x,y)
% \)
% since for any random variables $X$ and $Y$,
% \(
%     \EE\sbr{
%         \min(X,Y)
%     }
%     \geq 0
%     \Rightarrow
%     \EE\sbr{X},
%     \EE\sbr{Y}
%     \geq 0
% \).
% \kkcomment{When you move this up, state that the other welfares Nash s.t IR or social welfare s.t IR also work.}
% \mycomment{not necessarily true}

\subsection{Impossibility Under Strategic Reporting}

Thus far we have assumed that the agents report their valuations truthfully or that they can be publicly 
verified. This assumption is necessary to design a performant exchange algorithm; otherwise, competitors 
could misreport their valuations to manipulate the algorithm, increasing their own utility at the expense of overall welfare.
The following result demonstrates that the gap between the optimal algorithm when competitors are truthful and the optimal
algorithm when competitors are untruthful can be arbitrarily large.

\begin{restatable}{proposition}{optTruthfulVsOptUntruthful}
    \label{prop:opt-truthful-vs-opt-untruthful}
    For any $k\geq 0$ there exists an environment such that  
    \(
        \EE\sbr{
            \rewardss{
                \optstrategytt{
                    \optstoptime
                }
            }
        }
        -
        \EE\sbr{
            \rewardss{
                \opttruthstrategytt{
                    \opttruthstoptime
                }
            }
        }
        \geq k
    \),
    where 
    \(
        \optexpolicy
        =
        \rbr{
            \cbr{
                \optstrategyt
            }_{\timeind\geq 0},
            \optstoptime
        }
    \)
    and 
    \(
        \opttruthexpolicy
        =
        \rbr{
            \cbr{
            \opttruthstrategyt
            }_{\timeind\geq 0},
            \opttruthstoptime
        }
    \)
    are the optimal policies assuming and not assuming truthful reporting respectively.
    % \kkcomment{This should be stated more clearly.
    % You have defined $\optexpolicy$
    % and 
    % $\opttruthexpolicy$, but not used them. On the other hand, you haven't defined $\EE\sbr{
    %         \rewardss{
    %             \opttruthstrategytt{
    %                 \opttruthstoptime
    %             }
    %         }
    %     }$. }
    % \mycomment{put =}
\end{restatable}  
\begin{proof}
    Let $\numdataa=\numdatab=1$.
    Consider the valuations
    \begin{align*}
        \valaii{0}=\valbjj{0}=0
        \qquad
        \valaii{1}
        =
        \begin{cases}
            4k
            &
            \text{w.p. } \frac{1}{2}
            \\
            % \frac{k}{2}
            2k
            &
            \text{w.p. } \frac{1}{2}
        \end{cases}  
        \qquad
        \valbjj{1} = 4k
    \end{align*}  
    and let the competitors' utilities and the welfare function be given by
    \begin{align*}
        \utila(x,y) 
        =
        x-y
        \qquad
        \utilb(x,y)
        =
        \begin{cases}
            0 
            &
            \text{if }
            (y-x)\geq -2k
            %  k\geq x-y
            \\
            2(y-x)
            &
            \text{otherwise}
        \end{cases}  
        \qquad
        \welfarefunc(x,y)=x+y .
    \end{align*}
    First consider when competitors are truthful. 
    Under this model, the optimal policy 
    \(
        \optexpolicy
    \)
    first gives $\competitora$ competitor $\competitorb$'s point,
    if 
    $\valaii{1}=2k$ 
    then sharing terminates, otherwise $\competitorb$ is given
    $\competitora$'s point to ensure $\competitorb$'s utility does not lead to a negative welfare.
    Thus the expected welfare under this policy is 
    \[
        \EE\sbr{
            \rewardss{
                \stopstateforpolicy{
                    \optexpolicy_k
                }
            }
        }
        =
        \frac{1}{2}
        \welfarefunc\rbr{
            \utilass{
                (1,0)
            },
            \utilbss{
                (1,0)
            }
        }
        +
        \frac{1}{2}
        \welfarefunc\rbr{
            \utilass{
                (1,1)
            },
            \utilbss{
                (1,1)
            }
        }
        =
        \frac{1}{2}
        \rbr{
            % \frac{k}{2}
            2k
            + 0 
        }
        +
        \frac{1}{2}
        \rbr{
            0 + 0
        } 
        % =\frac{k}{4}
        =
        k.
    \] 
    Now consider when competitors are untruthful. 
    In this setting, the optimal truthful policy will
    not always 
    give $\competitorb$ competitor
    $\competitora$'s point. This is because doing so will always lead to a welfare of no more than 0,
    thus such a policy is no better than one which exchanges zero resources.
    Now consider a policy which only gives $\competitorb$ the point based on what is reported as $\valaii{1}$.
    In this case $\competitora$ is incentivized to always report the value for $\valaii{1}$ for which $\competitorb$
    does not receive the point as it increases their utility. However, this leads to a negative welfare.

    Therefore, the optimal policy when competitors are untruthul recommends no exchange so 
    \(
        \EE\sbr{
            \rewardss{
                \stopstateforpolicy{
                    \opttruthexpolicy_k
                }
            }
        }
        =0
    \)
    and thus
    \(
        \EE\sbr{
            \rewardss{
                \optstrategytt{
                    \optstoptime
                }
            }
        }
        -
        \EE\sbr{
            \rewardss{
                \opttruthstrategytt{
                    \opttruthstoptime
                }
            }
        }
        =k
    \).
\end{proof}  

For any $k\geq 0$ the proof of Proposition~\ref{prop:opt-truthful-vs-opt-untruthful} constructs an environment
where $\numdataa=\numdatab=1$ and $\valaii{1}$ is random but $\valbjj{1}$ is deterministic. When 
competitors are truthful, an optimal algorithm
relies on learning the realization of $\valaii{1}$ to determine whether to give $\competitorb$ competitor $\competitora$'s
resource. However, when competitors may be untruthful,
$\competitora$ is always incentivized to report whichever value leads to $\competitorb$ not obtaining
the resource. In this case the optimal exchange algorithm is to not recommend any exchanges.